\documentclass[aspectratio=169, 11pt]{beamer}
\usepackage[utf8]{inputenc}
\usepackage[T1]{fontenc}
\usepackage[french]{babel}
\usepackage{graphicx}
\usepackage{tikz}
\usepackage{xcolor}
\usepackage{hyperref}

\usetheme{Madrid}
\usecolortheme{beaver}

\setbeamercolor{frametitle}{fg=black,bg=white}
\setbeamercolor{title in head/foot}{fg=black,bg=white}

\setbeamertemplate{footline}{%
  \leavevmode%
  \hbox{%
    \begin{beamercolorbox}[wd=.82\paperwidth,ht=2.25ex,dp=1ex,leftskip=1.2ex]{author in head/foot}%
      \usebeamerfont{author in head/foot}%
      \footnotesize RVP2 --- Moteur de recherche multimodal (Ping) --- Barbier / Ledoux / Pascal / Wone
    \end{beamercolorbox}%
    \begin{beamercolorbox}[wd=.18\paperwidth,ht=2.25ex,dp=1ex,right,rightskip=1.2ex]{page number in head/foot}%
      \usebeamerfont{page number in head/foot}%
      \footnotesize \insertframenumber/\inserttotalframenumber
    \end{beamercolorbox}%
  }%
}

\title[LLM Filters]{Illustration du mode LLM Filters}
\author{Nathan Barbier - Projet Informatique}
\date{\today}

\begin{document}

\begin{frame}{LLM Filters : de l'intention utilisateur aux filtres}
    \begin{columns}[T,onlytextwidth]
        \begin{column}{0.58\textwidth}
            \centering
            % Remplacer par le nom du screenshot exporte
            \includegraphics[width=\linewidth,height=0.74\textheight,keepaspectratio]{llm_filters_slide.png}
        \end{column}

        \begin{column}{0.39\textwidth}
            {\small
            \begin{block}{Pourquoi c'est utile}
                Ce mode sert surtout a \textbf{reduire la friction} :
                l'utilisateur n'a pas besoin de connaitre tous les champs
                disponibles dans le moteur de recherche.
            \end{block}

            \begin{block}{Ce que montre cette capture}
                \begin{itemize}
                    \item une phrase libre est interpretee automatiquement
                    \item les elements importants sont convertis en filtres exploitables
                    \item on garde ensuite une recherche \textbf{fiable et verifiable}
                \end{itemize}
            \end{block}

            \begin{alertblock}{Point important}
                Le LLM n'est pas utilise pour classer directement les points.
                Il sert d'\textbf{interface de traduction} entre le langage naturel
                et les filtres deterministes du backend.
            \end{alertblock}
            }
        \end{column}
    \end{columns}
\end{frame}

\end{document}
